\subsection{Методы решения систем уравнений}
Будем рассматривать систему 
\begin{equation}
\label{eq:sae_generic}
\mat A u = f.
\end{equation}

\subsubsection{Простые итерационные алгоритмы}
\label{sec:sae_simple}
\subsubsubsection{Метод Якоби}
\label{sec:sae_jacobi}
Будем рассматривать систему уравнений вида
\begin{equation*}
    \sum_{j=0}^{N-1} A_{ij} u_j = r_i, \quad i = \overline{0, N-1}
\end{equation*}
относительно неизвестного сеточного вектора $\{u\}$.

В классическом виде алгоритм Якоби формулируется в виде
\begin{equation*}
    \hat u_i = \frac{1}{A_{ii}}\left(r_i - \sum_{j\neq i} A_{ij}{u_j}\right)
\end{equation*}

Произведём некоторые преобразования
\begin{align*}
    \hat u_i &= \frac{1}{A_{ii}}\left(r_i - \sum_{j} A_{ij}{u_j} + A_{ii}u_i\right) \\
             &= u_i + \frac{1}{A_{ii}}\left(r_i - \sum_{j} A_{ij}{u_j}\right)
\end{align*}

Таким образом, программировать итерацию этого алгоритма, обновляющую значения массива $\{u\}$, можно в виде
\begin{align*}
    &\check u = u; \\
    &\textbf{for } i=\overline{0, N-1} \\ 
    &\quad u_i \mathrel{{+}{=}} \frac{1}{A_{ii}}\left(r_i - \sum_{j=0}^{N-1} A_{ij}{\check u_j}\right)\\
    &\textbf{endfor} \\
\end{align*}

\subsubsubsection{Метод Зейделя}
\label{sec:sae_seidel}

Формулируется в виде
\begin{equation*}
    \hat u_i = \frac{1}{A_{ii}}\left(r_i - \sum_{j<i} A_{ij}{\hat u_j} - \sum_{j>i} A_{ij}{u_j} \right).
\end{equation*}

Поскольку этот метод неявный относительно уже найденных на итерации значений, то в отличии от метода Якоби этот алгоритм не требует создания временного массива $\hat u$
при программировании. Псевдокод для реализации итерации этого метода можно записать как

\begin{align*}
    &\textbf{for } i=\overline{0, N-1} \\ 
    &\quad u_i \mathrel{{+}{=}} \frac{1}{A_{ii}}\left(r_i - \sum_{j=0}^{N-1} A_{ij}{u_j}\right)\\
    &\textbf{endfor}
\end{align*}


\subsubsubsection{Метод последовательных верхних релаксаций SOR}
\label{sec:sae_sor}

Этот метод основан на добавлении к решению результатов итераций Зейделя с
коэффициентом $\omega > 1$. То есть он изменияет решение по тому же
принципу, что и метод Зейделя, но искуственно увеличивает эту добавку.

Формулируется этот метод в виде
\begin{equation*}
    \hat u_i = (1-\omega) u_i + \frac{\omega}{A_{ii}}\left(r_i - \sum_{j<i} A_{ij}{\hat u_j} - \sum_{j>i} A_{ij}{u_j} \right).
\end{equation*}
Для устойчивости метода необходимо $\omega < 2$. В частности, для одномерных задач,
заданных на единичном отрезке, для оптимальной сходимости можно использовать соотношение $\omega \approx 2 - 5 h$,
где $h$ -- шаг сетки.

Итерация этого метода по аналогии с методом Зейделя может быть запрограммирована в виде

\begin{align*}
    &\textbf{for } i=\overline{0, N-1} \\ 
    &\quad u_i \mathrel{{+}{=}} \frac{\omega}{A_{ii}}\left(r_i - \sum_{j=0}^{N-1} A_{ij}{u_j}\right)\\
    &\textbf{endfor}
\end{align*}

\subsubsection{Метод коррекции поправки}
\label{sec:it_def_corr}
Пусть $u^n$ -- некоторое приближённое решение системы \cref{eq:sae_generic} на итерационном слое.
Введём невязки как
\begin{equation}
\label{eq:sae_rn}
r^n = f - \left(\mat A u\right)^n
\end{equation}
и поправку для перехода на следующий итерационный слой:
\begin{equation*}
\triangle u^n = u^{n+1} - u^{n}.
\end{equation*}
Обобщённый процесс с предобуславливателем $\mat B$ для исходной системы можно записать в виде
\begin{equation}
\label{eq:sae_un_problem}
\mat B \triangle u^n = r^n
\end{equation}

Этот процесс является согласованным: если в качестве $u^n$ взять точное решение системы \cref{eq:sae_generic},
то $r^n = 0$, и следовательно, $\triangle u^n = 0$.

Алгоритм решения исходной системы будет иметь следующий вид:
\begin{enumerate}
\item задать начальное приближение $u^n$ при $n=0$
\item вычислить невязку $r^n$ из \cref{eq:sae_rn}
\item если норма невязки $||r^n|| < \epsilon$, то завершить итерации
\item решив систему \cref{eq:sae_un_problem} получить поправку $\triangle u^n = \mat B^{-1} r^n$.
\item найти следующее приближение $u^{n+1} = u^{n} + \triangle u^{n}$
\item увеличить счётчик $n = n + 1$ и перейти на шаг 2.
\end{enumerate}

Конкретный итерационный метод целиком определяется выбором предобуславливаетеля.
В качестве предобуславливателя обычно выбирают матрицу, похожую на $\mat A$, но при этом легкообратимую.
Некоторые примеры:
\begin{itemize}
\item если выбрать константу $B = 1/\gamma$, то получим
простой итерационный процесс Ричардсона с шагом $\gamma$.
\item диагональная часть исходного оператора $B = \{a_{ii}\}$ приводит к методу Якоби \secref{sec:sae_jacobi}
\item нижняя треугольная часть исходной матрицы $B = \{a_{ij}\}, j\leq i$ даёт метод Зейделя \secref{sec:sae_seidel}.
\end{itemize}

Если в качестве предобуславливателя взять изходную матрицу $\mat A$, 
то итерационный процесс сойдётся за одну итерацию. Действительно из \cref{eq:sae_rn,eq:sae_un_problem}:
\begin{equation*}
u^{n+1} = u^{n} + A^{-1}(f - A u^n) = A^{-1} f
\end{equation*}

Это метод требует прямого обращения только предобуславливаетеля $\mat B$, но не 
исходного оператора $\mat A$, поэтому он может быть применён в том числе и для решения нелинейных систем.

Для обращения самого предобуславливателя в свою очередь так же может
быть использован итерационный алгоритм. Этот итерационный процесс будет являться
внутриннем, поэтому не будет требовать глубокой сходимости.
Зачастую для обращения предобуславливателя достаточно сделать 2-3 поверхностные итерации
одним из простых методов из \secref{sec:sae_simple}.

\subsubsection{Методы подпространств Крылова}

\subsubsubsection{Итерационный процесс BiCGStab с предобуславливателем}
\label{sec:bicgstab}

Инициализация
\begin{enumerate}
\item есть начальное приближение $u^0$
\item вычисляем невязку
$$r_0 = f - \mat A u_0$$
\item предобуславливаем невязку
$$r_0' = \matinv{M}r_0$$
\item запоминаем начальную невязку
$$\tilde r = r_0'$$
\item инициализируем
$$p_0' = v_0' = 0, \quad \rho_0 = \omega_0 = \alpha_0 = 1.0$$
\end{enumerate}

Далее на $k$-ой итерации

\begin{enumerate}
\item проецируем текущую невязку на начальную
$$\rho_k = \tilde r^T \cdot r'_{k-1}$$
\item коэффициент $\beta$:
$$\beta_k = \frac{\rho_k}{\rho_{k-1}} \frac{\alpha_{k-1}}{\omega_{k-1}}$$
\item Обновляем направление
$$p'_k = r'_{k-1} + \beta_k \left( p'_{k-1} - \omega_{k-1} v'_{k-1} \right)$$
\item
$$v_k'=\matinv{M}\mat A p_k'$$
\item коэффициент $\alpha$
$$\alpha_k = \frac{\rho_k}{\tilde r^T \cdot v'_k}$$
\item Промежуточный вектор $s$:
$$s'_k =  r'_{k-1} - \alpha_k v'_k$$
\item Промежуточный вектор $t$
$$t'_k = \matinv{M}\mat{A}s_k'$$
\item множитель $\omega$
$$\omega_k = \frac{ (t_k')^T \cdot s'_k} { (t_k')^T \cdot t_k' }.$$
\item обновляем невязку и решение
\begin{equation}
\label{eq:schwarz_bicgstab}
\begin{aligned}
& u_k = u_{k-1} + \alpha_k p'_k + \omega_k s'_k, \\
& r'_k = s'_k - \omega_k t'_k
\end{aligned}
\end{equation}
\end{enumerate}
